\section{Identified Gap and Positioning of the Thesis}
\label{sec:gaps}

The preceding sections have described a mature research landscape. Cooperative game theory provides the canonical reference for benefit allocation \cite{shapley1953, moncecchi2020}; a set of scalable approximations has made its computation feasible also for large communities \cite{castro2009, cremers2023, alonsopedrero2024, ferrucci2025}; simplified approaches remain widely used because of their transparency towards non-expert users \cite{kulmala2021, fioriti2023}; and a group of fairness indicators, imported from other disciplines, is available to evaluate the resulting allocations (Section~\ref{sec:2.3}). The open problem is therefore no longer how to \emph{compute} a fair allocation, but how to \emph{choose} among allocations that can all be computed.

\subsection*{The research need from which this work originates}

The most extensive attempt to answer this question is the data-driven comparison of González-Asenjo and Izquierdo (2026), which evaluates five methods belonging to three methodological families over almost fourteen thousand simulated configurations, varying community size up to one thousand members, demand heterogeneity, renewable penetration and the internal pricing regime \cite{gonzalezAsenjo2026, schmeidler1969}. Their main result is that no method dominates across all configurations and that the choice of the allocation mechanism remains dependent on the operating context.

When discussing the limitations of this result, the authors explicitly indicate the direction in which it should be extended. The community archetype they adopt is deliberately kept \textit{``sufficiently abstract to allow for systematic comparison across different allocation rules and operating conditions''}, with a generic internal market replacing any real remuneration scheme, and the analysis proceeds by \textit{``abstracting from investment decisions, network constraints, and subjective behavioural factors''}. Their conclusion is explicit: \textit{``regulatory constraints and tariff structures may further influence feasible allocation schemes in practice. Addressing these aspects represents a natural direction for future research''} \cite{gonzalezAsenjo2026}.

This research need is not a marginal one, and its relevance can be assessed by considering what the same authors identify as the driver of their own result. In their model the value to be distributed is the surplus generated by internal trading, that is, the area between the willingness to pay of net consumers and the selling prices of net producers. Consistently, they observe that \textit{``price heterogeneity emerges as a central driver of CAMs' behaviour''} and they warn that \textit{``conclusions drawn under uniform pricing may not generalize to heterogeneous market contexts''} \cite{gonzalezAsenjo2026}. There are, however, regimes in which this driver is absent by construction. When the remuneration of energy sharing is a regulated transfer rather than an arbitrage, the shared quantity keeps the same structure, namely the hourly minimum between aggregate injections and aggregate withdrawals, but its valorization does not depend on individual prices, and the identity of the members who contribute does not change the amount to be distributed. Together with the internal price, the only endogenous criterion for ranking alternative allocation rules also disappears. Such regimes are not marginal, since most European energy communities are based on virtual metering, and this condition had already been noted by Minuto and Lanzini, who observed that \textit{``the lack of a deep and detailed investigation on the sharing mechanism in this energy policy framework results in a significant scientific literature gap''} \cite{minutoLanzini2022}. Furthermore, the formulation adopted by González-Asenjo and Izquierdo is stronger than it may appear at a first reading: regulatory constraints do not simply modify the outcome of the allocation, they influence which schemes are \emph{feasible}. In at least one operative framework, regulation sets a threshold beyond which a share of the remuneration cannot be distributed to a given class of members, so that any allocation exceeding that threshold is not admissible \cite{MASE:Decreto:CER:2023}.

The gap on which this work is based is therefore the one that those authors declare they have not addressed. It is not known how benefit allocation methods rank against each other when the value to be distributed is no longer a market surplus proportional to price dispersion but a regulated transfer, and when regulation itself restricts the set of admissible allocations. The reference comparison available in the literature has been obtained under the first condition, and whether its conclusions can be transferred to the second has not been verified.

\subsection*{Contribution of the thesis and its extensions}

This thesis addresses this research need by repeating the comparison under those conditions, by means of a techno-economic simulation model in which the amount available for allocation and the constraints acting on it are determined by a regulated remuneration scheme, and in which the allocation mechanisms are evaluated under identical operating conditions. Starting from this core, the work extends the analysis in three directions, each of which is also motivated by the literature.

The first direction concerns the set of mechanisms that are compared. The authors of the reference study observe that \textit{``the modular nature of this framework allows for the evaluation of methods not considered in this paper''} \cite{gonzalezAsenjo2026}, and their design covers three of the six families identified in Section~\ref{sec:2.2}. The extension to the remaining families is not motivated only by completeness. Once the internal price is removed, the temporal alignment between withdrawal and injection is the only dimension along which members differ, and this is precisely the dimension that performance-based mechanisms are designed to capture \cite{gianaroli2024, giudice2026}.

The second direction concerns the evaluation layer. The same authors conclude that \textit{``formal property compliance alone is insufficient to assess real-world performance and must be complemented by scenario-specific analysis''} \cite{gonzalezAsenjo2026}, a statement that connects the present discussion to the axiomatic framework presented in Section~\ref{sec:2.5} \cite{gonzalezAsenjo2025}. The instrument required for such an analysis is, however, not well established. Dynge and Cali observe that the fairness indicators currently in use were \textit{``adapted for use in the context of local electricity markets''} from other disciplines and that, before their study, \textit{``no research has been found addressing or comparing the performance of these indicators''}; they evaluate four of them and state that \textit{``additional investigation is needed to refine fairness indicators, ensuring their effectiveness across diverse market models''} \cite{dynge2025}. A similar need is expressed independently by Volpato et al., who indicate among future developments the identification of parameters able to compare benefit distributions in quantitative terms \cite{volpato2024}. For this reason the present work employs a set of indicators rather than a single one, and treats their agreement as a result to be observed.

The third direction concerns the configurations over which the comparison is repeated. Gasca et al.\ have already crossed mechanisms and indicators, evaluating ten approaches with four indices at the same time, and conclude that \textit{``existing fairness indexes do not adequately represent fairness''} and that \textit{``the findings are sensitive to user assets and profile changes''}. They also state, however, that \textit{``the same conclusions apply to other types of ECs with diverse user types''} \cite{gasca2024}. This is a claim of generality with respect to composition that their design does not verify, since the study is based on a single case study in which composition is not one of the factors that are varied. The fact that such a statement immediately follows the acknowledgement of sensitivity to user profiles shows the distance between what this literature assumes and what it has actually verified. The same need appears in different forms in the studies that indicate as future developments \textit{``the different EC compositions (number of members in the UT groups)''} \cite{minutoLanzini2022}, \textit{``the critical set-up''} beyond which a mechanism is no longer sensibly applicable \cite{fina2024}, and \textit{``the number and composition of community members''} \cite{giudice2026}. The comparison is therefore repeated over configurations that systematically vary the number of members, the prosumer penetration and the composition by member category.

\subsection*{Research question and empirical context}

The gap and its extensions converge into a single research question:

\begin{quote}
\emph{When the value to be distributed within an energy community is a regulated transfer rather than a market surplus, and regulation itself restricts the set of admissible allocations, how do allocation mechanisms rank against each other in terms of distributive equity, incentive strength of the rule and economic sustainability of participation; to what extent do different fairness indicators agree in evaluating them; and to what extent does this ranking depend on the composition of the community?}
\end{quote}

The analysis is carried out within the Italian regulatory framework, which represents the empirical context of the work and not its motivation. This framework is an operative and fully documented instance of the regulated remuneration regime described above, and it is examined in Chapter~\ref{ch:italian_context}. The legally defined member categories that determine the remuneration treatment within it enter the model as constraints of the simulated framework.

The thesis answers the research question through a simulation model organised in two layers. The valorisation layer has no degrees of freedom: starting from the hourly profiles of the members, it determines the amount available for allocation and the constraints acting on it. The allocation layer is the object of the study, and implements one representative mechanism for each of the families of Section~\ref{sec:2.2}, selected according to the axiomatic criterion of Section~\ref{sec:2.5}. The operational definition of the indicators and the characterisation of the underlying game are developed in the methodological chapter. The expected outcome is not the identification of an optimal mechanism, but the structure of the interaction between the three elements: the configurations in which the ranking of the mechanisms is reversed, and those in which the indicators agree or diverge.
